\subsection{Documentos a Utilizar en el Proyecto de Software}

\subsubsection*{D-01 Acta de Constitución del Proyecto}
\textbf{Propósito:} Autorizar formalmente el inicio del proyecto y designar al Project Manager.  
\textbf{Responsable:} Patrocinador del Proyecto.  
\textbf{Fase:} Inicio.

\subsubsection*{D-02 Registro de Stakeholders}
\textbf{Propósito:} Identificar, clasificar y analizar a los interesados del proyecto, definiendo su nivel de influencia y participación.  
\textbf{Responsable:} Project Manager.  
\textbf{Fase:} Inicio y actualización durante la planificación.

\subsubsection*{D-03 Plan de Gestión del Alcance}
\textbf{Propósito:} Definir cómo se gestionará, validará y controlará el alcance del proyecto.  
Incluye el Scope Statement, la EDT/WBS, la línea base del alcance y el procedimiento de control de cambios.  
\textbf{Responsable:} Project Manager.  
\textbf{Fase:} Planificación.

\subsubsection*{D-04 Documento de Requerimientos (SRS)}
\textbf{Propósito:} Especificar los requisitos funcionales y no funcionales del sistema, incluyendo interoperabilidad, seguridad de datos y requisitos de desempeño.  
\textbf{Responsable:} Analista Funcional.  
\textbf{Fase:} Planificación y refinamiento iterativo.

\subsubsection*{D-05 Matriz de Trazabilidad de Requisitos}
\textbf{Propósito:} Asegurar la relación directa entre requisitos, componentes desarrollados y casos de prueba.  
\textbf{Responsable:} Analista Funcional / Líder QA.  
\textbf{Fase:} Ejecución y Monitoreo.

\subsubsection*{D-06 Documento de Arquitectura del Software}
\textbf{Propósito:} Definir la arquitectura tecnológica del sistema, modelo cloud, integración con APIs externas y mecanismos de seguridad.  
\textbf{Responsable:} Arquitecto de Software.  
\textbf{Fase:} Planificación técnica y Ejecución.

\subsubsection*{D-07 Documento de Diseño Técnico}
Incluye diagramas UML, modelo entidad–relación y diagramas de integración.  
\textbf{Responsable:} Equipo de Arquitectura.  
\textbf{Fase:} Planificación y Ejecución.

\subsubsection*{D-08 Plan de Iteraciones}
\textbf{Propósito:} Definir las funcionalidades a desarrollar en cada iteración bajo el enfoque iterativo e incremental.  
\textbf{Responsable:} Project Manager / Equipo Técnico.  
\textbf{Fase:} Ejecución.

\subsubsection*{D-09 Registro de Control de Cambios}
\textbf{Propósito:} Documentar y gestionar las solicitudes de cambio que impacten el alcance, cronograma o costo.  
\textbf{Responsable:} Project Manager.  
\textbf{Fase:} Monitoreo y Control.

\subsubsection*{D-10 Plan de Pruebas}
Incluye pruebas funcionales, de integración, seguridad y desempeño.  
\textbf{Responsable:} Líder de Aseguramiento de Calidad (QA).  
\textbf{Fase:} Ejecución.

\subsubsection*{D-11 Informe de Pruebas y Validación}
\textbf{Propósito:} Evidenciar el cumplimiento de los criterios de aceptación definidos para cada entregable.  
\textbf{Responsable:} Equipo QA / Analista Funcional.  
\textbf{Fase:} Monitoreo y Control.

\subsubsection*{D-12 Manual Técnico}
\textbf{Propósito:} Documentar instalación, configuración y mantenimiento del sistema.  
\textbf{Responsable:} Arquitecto de Software.  
\textbf{Fase:} Cierre.

\subsubsection*{D-13 Manual de Usuario}
\textbf{Propósito:} Guiar al personal administrativo en el uso adecuado del sistema.  
\textbf{Responsable:} Equipo de Proyecto.  
\textbf{Fase:} Transición.

\subsubsection*{D-14 Plan de Capacitación}
\textbf{Propósito:} Definir cronograma y contenidos de formación para el personal.  
\textbf{Responsable:} Project Manager / Analista Funcional.  
\textbf{Fase:} Transición.

\subsubsection*{D-15 Actas de Aceptación Parcial}
\textbf{Propósito:} Formalizar la validación incremental de los entregables por iteración.  
\textbf{Responsable:} Representante del Cliente.  
\textbf{Fase:} Ejecución y Control.

\subsubsection*{D-16 Acta de Cierre del Proyecto}
\textbf{Propósito:} Formalizar la aceptación final del sistema y el cierre administrativo del proyecto.  
\textbf{Responsable:} Patrocinador del Proyecto.  
\textbf{Fase:} Cierre.